\documentclass[11pt,a4paper]{article}
\usepackage[utf8]{inputenc}
\usepackage{geometry}
\geometry{top=1in, bottom=1in, left=1in, right=1in}
\usepackage{amsmath, amssymb, amsfonts, bm}
\usepackage{graphicx}
\usepackage{booktabs}
\usepackage{hyperref}
\usepackage{tikz}
\usetikzlibrary{positioning, calc, arrows.meta}

% Color definitions
\definecolor{primary}{RGB}{30, 41, 59}       % Dark Slate
\definecolor{secondary}{RGB}{14, 116, 144}   % Teal Accent

\hypersetup{
    colorlinks=true,
    linkcolor=secondary,
    citecolor=secondary,
    urlcolor=secondary
}

\title{\textbf{Compiling Online Reasoning into Offline Learned Sparse Visual Indexes:\\A HIVE-to-V-SPLADE Distillation Framework}}
\author{\textbf{Esteban Schafir} \\ FIU Research}
\date{\today}

\begin{document}

\maketitle

\begin{abstract}
Multimodal Information Retrieval (MMIR) is currently divided into two suboptimal paradigms: low-latency learned sparse retrievers that lack complex reasoning capability, and high-latency LLM-driven pipelines that resolve reasoning gaps online at high computational cost. In this proposal, we introduce a teacher-student distillation framework that compiles the multi-step online reasoning loop of a hypothesis-driven retriever (HIVE) into a static, query-encoding-free visual sparse index (V-SPLADE) offline. By supervising the student sparse encoder with reasoning-aware gating targets and listwise rank margins, we bake query-time reasoning directly into visual document term expansions. This architecture allows complex reasoning visual search to be executed in sub-10ms using standard inverted index lookup engines.
\end{abstract}

---

\section{Introduction and Theoretical Motivation}
Modern Visual Document Retrieval (VDR) requires matching visual pages (e.g., charts, schematics, academic papers) to text queries. Historically, this has been solved using two extremes:
\begin{enumerate}
    \item \textbf{Learned Sparse Retrieval (LSR):} Models like SPLADE and V-SPLADE map document images offline to a vocabulary-indexed sparse space $w_p \in \mathbb{R}^{|V|}$. At search time, a text query $q$ is matched via simple dot-product calculations ($w_q^T w_p$) over inverted indexes, yielding sub-10ms latency. However, these models cannot handle complex queries that require multi-hop visual reasoning because they are trained using simple, literal image captions.
    \item \textbf{Reasoning-Intensive Pipelines:} Frameworks like HIVE deploy Large Language Models (LLMs) at query time to analyze initial search candidate documents, synthesize compensatory queries, and verify/rerank the final list. While achieving state-of-the-art accuracy on reasoning-intensive benchmarks (e.g., MM-BRIGHT), they require multiple LLM inferences per query, adding seconds of latency and massive API costs.
\end{enumerate}

The core research question of this proposal is: \textit{Can we compile the reasoning capability of a multi-stage online LLM pipeline into a static visual sparse representation offline?} 

We propose a teacher-student framework to perform this compilation. Under this framework, the online HIVE pipeline serves as the \textbf{Teacher} during training, producing reasoning targets. The V-SPLADE model serves as the \textbf{Student}, learning to pre-activate reasoning-related tokens directly from image pixels so that no LLM or deep vision network is needed at query time.

---

\section{Formulation of the Distillation Framework}

The architecture of the proposed distillation pipeline is illustrated in Figure~\ref{fig:distill_arch}.

\begin{figure}[htbp]
\centering
\includegraphics[width=\textwidth]{figures/Pipeline.png}
\caption{Architecture of the complete HIVE-to-V-SPLADE Distillation Pipeline, showcasing: (1) Offline Teacher Generation utilizing local or NIM API LLM reranking; (2) Student Distillation focusing exclusively on optimizing the Query LUT weights; and (3) Online serving via a Hybrid Inverted/Vector DB index combining sparse (V-SPLADE) and dense (BGE/FAISS) retrieval.}
\label{fig:distill_arch}
\end{figure}

\subsection{The Teacher Model (HIVE)}
Given a dataset of training queries $Q$ and a visual document corpus $C$, HIVE acts as the teacher. For a given training query $q = (q_t, q_v)$ and document image $d \in C$, HIVE generates:
\begin{itemize}
    \item \textbf{A listwise verification score $S_{\text{HIVE}}(d \mid q)$:} Representing the reasoning-based relevance of document $d$ to query $q$.
    \item \textbf{A compensatory query $\hat{q}$:} A sequence of tokens synthesized by the LLM (Stage 2) representing the logical concepts missing from initial candidate pools.
    \item \textbf{A verification rationale $R(d, q)$:} The step-by-step reasoning explanation generated by the LLM (Stage 4) verifying why document $d$ matches query $q$.
\end{itemize}

\subsection{The Student Model (V-SPLADE)}
The student is a visual learned sparse encoder. It takes the document image $d$ and maps its visual tokens to a sparse representation $w_p \in \mathbb{R}^{|V|}$:
\begin{equation}
w_p[v] = \max_{t=1 \dots L} \log(1 + \text{ReLU}(\text{LMHead}(h_t)[v]))
\end{equation}
where $h_t$ are visual hidden states and $v \in V$ is a vocabulary token.

\subsection{Distillation Objectives}
To transfer reasoning capabilities, we introduce two distinct training objectives:

\subsubsection{Objective 1: Reasoning-Aware Gated Supervision (Lexical Distillation)}
In standard V-SPLADE, vocabulary activations are gated by literal caption tokens. To encode reasoning, we compile a target reasoning vocabulary vector $W_R \in \mathbb{R}^{|V|}$ by tokenizing the union of the original query $q_t$, the LLM-synthesized compensatory query $\hat{q}$, and key terms in the verification rationale $R(d, q)$:
\begin{equation}
W_R = \text{BagOfWords}(q_t \cup \hat{q} \cup \text{Keywords}(R(d, q)))
\end{equation}

We define the overlap gate $o_{\text{reason}}[v]$ for the student's sparse document vector as:
\begin{equation}
o_{\text{reason}}[v] = w_p^{\text{sg}}[v] \cdot W_R[v]
\end{equation}
where $w_p^{\text{sg}}$ is the stop-gradient representation of the image logits. We sharpen and normalize this gate to compute token targets $\alpha_{\text{reason}}[v]$:
\begin{equation}
\alpha_{\text{reason}}[v] = \frac{o_{\text{reason}}[v]^{1/\tau}}{\sum_{v'} o_{\text{reason}}[v']^{1/\tau}}
\end{equation}
The student is supervised using the cross-entropy loss against these reasoning-expanded targets:
\begin{equation}
\mathcal{L}_{\text{lexical}} = -\frac{1}{B} \sum_{i=1}^B \sum_v \alpha_{\text{reason}, i}[v] \log \sigma(z_{p_i}[v])
\end{equation}
where $z_p[v]$ are the visual logits before activation.

\subsubsection{Objective 2: Rank Alignment Distillation (Margin MSE)}
To match the ranking outputs of HIVE, we minimize the difference between the student's inner product scores and the teacher's listwise reranking scores over document candidate pairs $(d_i, d_j)$:
\begin{equation}
\mathcal{L}_{\text{rank}} = \sum_{(d_i, d_j)} \left| (s_{\text{SPLADE}}(q, d_i) - s_{\text{SPLADE}}(q, d_j)) - (S_{\text{HIVE}}(d_i \mid q) - S_{\text{HIVE}}(d_j \mid q)) \right|^2
\end{equation}
where $s_{\text{SPLADE}}(q, d) = w_q^T w_p$ is V-SPLADE's query-indexing matching score.

The total student loss is:
\begin{equation}
\mathcal{L}_{\text{total}} = \mathcal{L}_{\text{lexical}} + \lambda_1 \mathcal{L}_{\text{rank}} + \lambda_2 \mathcal{L}_{\text{sparsity}}
\end{equation}
where $\mathcal{L}_{\text{sparsity}}$ is the $L1$ regularization on document weights to keep the index sparse.

---

\section{Detailed Use Case Walkthroughs}

\subsection{Example 1: Electrical Circuit Analysis (Logical Visual Dependencies)}
Consider a visual document database containing circuit diagrams.
\begin{itemize}
    \item \textbf{The Document Image $d$ (assigned as Doc 25):} A diagram showing a battery, a switch in the \textbf{open position} (depicted as a visual gap in the wiring), and a light-emitting diode (LED).
    \item \textbf{The Query $q$:} Text query \texttt{"Why is my LED not lighting up?"}.
    \item \textbf{The HIVE Teacher Loop (Offline Training):}
    \begin{enumerate}
        \item \textit{Initial Retrieval:} Finds documents discussing dead batteries or burnt-out resistors.
        \item \textit{Query Synthesis:} The teacher LLM analyzes the layout of Doc 25, detects the open switch, and generates the compensatory query $\hat{q}$: \texttt{"open switch loop circuit current flow LED"}.
        \item \textit{Verification:} The LLM generates the rationale $R(\text{Doc 25}, q)$: \textit{"The LED fails to illuminate because the diagram depicts an open switch, which breaks circuit loop continuity, preventing current flow."}
    \end{enumerate}
    \item \textbf{Reasoning Vocabulary Extraction:} The gating vector $W_R$ is compiled from the key terms: \{\texttt{switch}, \texttt{open}, \texttt{continuity}, \texttt{break}, \texttt{current}, \texttt{flow}\}.
    \item \textbf{Student Mapping Distillation:} The student V-SPLADE visual encoder learns to map the specific visual features of the wire gap in Doc 25 directly to these logical reasoning tokens during training.
    
    \item \textbf{Offline Indexing Phase (Static Database Construction):} 
    When indexing the corpus, the fine-tuned, frozen V-SPLADE model runs once over Doc 25. It detects the open switch and outputs the following sparse vector:
    \begin{equation}
    w_{\text{Doc 25}} = \{\text{\texttt{circuit}}: 4.5, \text{\texttt{diagram}}: 3.8, \text{\texttt{open}}: 4.0, \text{\texttt{break}}: 3.5, \text{\texttt{continuity}}: 3.0\}
    \end{equation}
    These weights are saved in a static \textbf{inverted index} (the search engine's internal dictionary):
    \begin{align*}
        \text{\texttt{circuit}} &\to [(\text{Doc 1}, 2.1), (\text{Doc 25}, 4.5)] \\
        \text{\texttt{open}} &\to [(\text{Doc 25}, 4.0)] \\
        \text{\texttt{break}} &\to [(\text{Doc 25}, 3.5), (\text{Doc 88}, 1.2)] \\
        \text{\texttt{continuity}} &\to [(\text{Doc 25}, 3.0)]
    \end{align*}
    
    \item \textbf{Query-Time Inference (Sub-10ms Serving):} 
    When a user searches for \texttt{"Why is my circuit broken?"}:
    \begin{enumerate}
        \item The text query is tokenized into \texttt{circuit} and \texttt{broken} (stemmed to \texttt{break}).
        \item The system performs a standard dictionary lookup on the inverted index under \texttt{circuit} and \texttt{break}.
        \item It immediately retrieves Doc 25 and computes the matching score via dot product: 
        \begin{equation}
        \text{Score} = w_q[\text{\texttt{circuit}}] \cdot 4.5 + w_q[\text{\texttt{break}}] \cdot 3.5 = 8.0
        \end{equation}
        This lookup is completely static, requires \textbf{no GPU execution or active LLMs at query time}, yet returns the correct document containing the visual reasoning explanation.
    \end{enumerate}
\end{itemize}

\subsection{Example 2: Financial Report Layouts (Data Chart Reasoning)}
Consider a visual database containing corporate quarterly reports.
\begin{itemize}
    \item \textbf{The Document Image $d$:} A page displaying a bar chart showing declining revenue over four consecutive quarters.
    \item \textbf{The Query $q$:} Text query \texttt{"Which companies show consecutive quarters of negative growth?"}
    \item \textbf{The HIVE Teacher Loop:}
    \begin{enumerate}
        \item \textit{Initial Retrieval:} Matches documents mentioning "negative growth" in the text, missing reports where the decline is only shown visually in the chart.
        \item \textit{Query Synthesis:} LLM analyzes the chart, notes the declining bars, and synthesizes $\hat{q}$: \texttt{"revenue decline chart consecutive quarters negative slope trend"}.
        \item \textit{Verification:} The LLM verifies the document with $R$: \textit{"The chart shows four declining bars representing successive quarterly revenue drops."}
    \end{enumerate}
    \item \textbf{Lexing Gating Target $W_R$:} Contains \{\texttt{chart}, \texttt{consecutive}, \texttt{decline}, \texttt{revenue}, \texttt{slope}, \texttt{drop}\}.
    \item \textbf{Student Distillation:} V-SPLADE's visual encoder learns to map the visual representation of descending bar charts to lexical tags like \texttt{decline}, \texttt{negative}, and \texttt{quarters}.
    \item \textbf{Query-Time Inference:} A text search matching "consecutive quarters of negative growth" hits the document because the visual encoder pre-activated those semantic concepts in the index offline.
\end{itemize}

\section{Empirical Implementation and Optimization}
We describe the current concrete implementation of the distillation framework, detailing the engineering optimizations and architectural decisions required for convergence and zero-shot out-of-domain (OOD) generalizability.

\subsection{Backbone and Parameterization}
The student model utilizes a frozen \texttt{Qwen/Qwen2.5-7B-Instruct} (or \texttt{Qwen/Qwen2.5-1.5B-Instruct}) backbone. The vocabulary size $|V|$ is $151,936$, and the hidden state dimensionality is $3584$ ($1536$ for 1.5B). Because fine-tuning the 544M parameter vocabulary projection head (\texttt{lm\_head}) on a small distillation dataset triggers severe overfitting and catastrophic forgetting of zero-shot vocabulary mappings on OOD splits, we parameterize the distillation as:
\begin{enumerate}
    \item \textbf{Frozen Document Representation (Path A):} The visual document encoder is kept 100\% frozen (including the pre-trained backbone and the projection head). This guarantees that the encoder's strong pre-trained linguistic and visual associations remain intact.
    \item \textbf{Trainable Query LUT:} A flat query term weight lookup table $L_{\text{LUT}} \in \mathbb{R}^{|V|}$ is optimized during training. It maps query token IDs directly to weights ($w_q = \text{softplus}(L_{\text{LUT}}[q])$), functioning as a vocabulary-level term importance map learned from the teacher's rationale patterns.
\end{enumerate}

\subsection{Memory-Efficient Chunked Logit Projection}
Projecting hidden states of shape $(B, L, 1536)$ directly into the $151,936$-dimensional vocabulary space consumes excessive GPU memory, leading to Out-of-Memory (OOM) failures on standard hardware. We mitigate this by calculating logits and performing max-pooling in chunks along the vocabulary dimension:
\begin{equation}
z_{p, c} = \max_{t} \left( H_t W_{\text{head}, c}^T \right)
\end{equation}
where $W_{\text{head}, c}$ is a sub-matrix of \texttt{lm\_head} of size $20,000 \times 1536$. The chunked max-pooled vectors are concatenated on the fly, reducing memory overhead by over 80\%.

\subsection{Query-Level Stop-Word Gating (BPE Compatibility)}
Applying stop-word vocabulary masking on Byte-Pair Encoding (BPE) subwords destroys content word components. For example, BPE tokenizes \texttt{"beavers"} into \texttt{["be", "avers"]}. Since \texttt{"be"} is a common English stop-word, a token-level mask zeroes it out, breaking the search key.
To resolve this, we implement \textbf{string-level query stopping}:
\begin{itemize}
    \item Stop-words and HTML/punctuation tags are removed from the query string \textit{prior} to tokenization.
    \item The document's raw tokenizations remain untouched, preserving BPE tokenization integrity while ensuring stop-word activations are ignored in query dot product calculations.
\end{itemize}

\subsection{Training and Evaluation Setup}
We train the model using:
\begin{itemize}
    \item \textbf{Optimizer:} AdamW optimizer applied exclusively to the Query LUT parameters ($L_{\text{LUT}}$) with a learning rate of $5 \times 10^{-3}$ and weight decay of $0.01$.
    \item \textbf{Loss Weights:} We configure $\lambda_{\text{rank}} = 1.0$, $\lambda_{\text{sparsity}} = 1 \times 10^{-4}$, and gating temperature $\tau = 1.0$.
    \item \textbf{Serving Latency:} Since document encoding is fully decoupled from query serving, matching is performed via simple CPU-level dot products, yielding a latency of under 5ms.
\end{itemize}

\subsection{Distillation Data Evolution: Multi-Phase Dataset Design}
To train the vocabulary lookup table $L_{\text{LUT}}$ (which maps Qwen's $151,936$-dimensional vocabulary directly to retrieval weights), we designed a multi-phase target data generation process, scaling the number of targets and expanding vocabulary domain diversity while maintaining a strict zero-shot held-out split protocol.

\subsubsection{Teacher Target Generation Architecture}
Teacher target generation is performed offline using high-capacity LLMs:
\begin{itemize}
    \item \textbf{Teacher Models:} Offline target generation utilizes the state-of-the-art \texttt{meta/llama-3.1-70b-instruct} model hosted via the NVIDIA NIM API, as well as local execution of \texttt{Qwen/Qwen2.5-32B-Instruct}.
    \item \textbf{Target Extraction:} For each query, HIVE Stage 2 generates a compensatory query $\hat{q}$, while Stage 4 generates step-by-step rationales $R(d, q)$ and continuous relevance scores $S_{\text{HIVE}}(d \mid q) \in [0.0, 5.0]$ across a 20-document candidate pool.
\end{itemize}

\subsubsection{Target Dataset Composition (5,055 Total Targets)}
To ensure robust zero-shot evaluation without data leakage, we explicitly reserved 9 target domains exclusively for held-out inference (\texttt{law}, \texttt{gaming}, \texttt{askubuntu}, \texttt{salesforce}, \texttt{travel}, \texttt{academia}, \texttt{medicalsciences}, \texttt{math}, \texttt{psychology}). All training targets were compiled strictly from non-overlapping splits, yielding \textbf{5,055 total distillation targets}:
\begin{enumerate}
    \item \textbf{MM-BRIGHT Non-Evaluation Training Splits (1,196 Targets):} Query-document targets compiled across 20 distinct scientific and technical MM-BRIGHT training splits (\texttt{aviation}: 125, \texttt{biology}: 99, \texttt{physics}: 95, \texttt{bioinformatics}: 89, \texttt{quantumcomputing}: 87, \texttt{earthscience}: 85, \texttt{crypto}: 74, \texttt{chemistry}: 65, \texttt{bitcoin}: 64, \texttt{sustainability}: 62, \texttt{philosophy}: 50, \texttt{pm}: 50, \texttt{gis}: 44, \texttt{bioacoustics}: 41, \texttt{quant}: 34, \texttt{economics}: 31, \texttt{christianity}: 30, \texttt{robotics}: 30, \texttt{islam}: 27, \texttt{apple}: 14).
    \item \textbf{Text-Rich BEIR Benchmarks (3,855 Targets):} To resolve out-of-vocabulary (OOV) bottlenecks across general and domain-specific vocabularies, targets were compiled from five BEIR corpora:
    \begin{itemize}
        \item \textbf{\texttt{BeIR/hotpotqa} (Multi-hop QA):} 1,567 queries providing multi-hop reasoning patterns across general knowledge.
        \item \textbf{\texttt{BeIR/scidocs} (Science/Bio):} 998 academic citation queries covering biomedical and physical abstracts.
        \item \textbf{\texttt{BeIR/fiqa} (Finance):} 500 queries providing corporate and financial terminology.
        \item \textbf{\texttt{BeIR/dbpedia-entity} (General Entity):} 467 structured entity queries (400 test + 67 dev) covering history, arts, and entity linking.
        \item \textbf{\texttt{BeIR/nfcorpus} (Clinical Nutrition):} 323 clinical medical and nutritional science queries.
    \end{itemize}
    \item \textbf{Synthetic Layout QA (4 Targets):} 4 query targets from \texttt{vidore/syntheticDocQA\_government\_reports\_test}.
\end{enumerate}

Table~\ref{tab:target_dataset_breakdown} details the complete target dataset breakdown across all source corpora.

\begin{table}[htbp]
\centering
\scriptsize
\caption{Composition of the Clean Distillation Target Dataset (5,055 Total Queries) strictly excluding the 9 held-out evaluation test splits.}
\label{tab:target_dataset_breakdown}
\begin{tabular}{lccl}
\toprule
\textbf{Source Dataset} & \textbf{Targets} & \textbf{Percentage} & \textbf{Domain / Semantic Focus} \\
\midrule
\texttt{beir/hotpotqa} & 1,567 & 31.0\% & Multi-hop Wikipedia reasoning \& general knowledge \\
\texttt{beir/scidocs} & 998 & 19.7\% & Scientific \& biomedical citation matching \\
\texttt{mm-bright/MM-BRIGHT} & 1,196 & 23.7\% & 20 multimodal scientific/technical training domains \\
\texttt{beir/fiqa} & 500 & 9.9\% & Corporate finance, banking, and economics QA \\
\texttt{beir/dbpedia-entity} & 467 & 9.2\% & Structured knowledge base entity retrieval \\
\texttt{beir/nfcorpus} & 323 & 6.4\% & Clinical nutrition and biomedical literature \\
\texttt{vidore/syntheticDocQA} & 4 & 0.1\% & Synthetic layout document QA \\
\midrule
\textbf{Total} & \textbf{5,055} & \textbf{100.0\%} & \textbf{Full vocabulary coverage ($|V|=151,936$)} \\
\bottomrule
\end{tabular}
\end{table}

\subsubsection{Why Visual Datasets Were Replaced in Training}
Although our student model is evaluated on visual document images in production, \textbf{the student's document encoder and Query LUT operate entirely in the text domain during training}. During indexing, the visual layouts are pre-parsed using frozen document layout models, mapping the visual content to textual BPE tokens.
Switching from synthetic visual datasets (like \texttt{syntheticDocQA}) to standard text-rich BEIR datasets provided two critical advantages:
\begin{itemize}
    \item \textbf{Continuous Relevance Gradients:} Synthetic visual datasets lack document text, providing only binary Q\&A pairs. This causes the teacher LLM to assign flat binary scores (e.g., $0.0$ or $5.0$). Text-rich BEIR datasets contain full-length documents, allowing the teacher to assign granular relevance scores.
    \item \textbf{Namespace Collision Mitigation:} Query IDs are namespaced (e.g., prefixed as \texttt{beir\_fiqa\_test\_q}) to prevent key conflicts with existing targets.
\end{itemize}

\section{Empirical Results and Analysis}
We evaluated the distilled V-SPLADE student model (incorporating Path A LUT-only training and pre-tokenization query stopping) across the 9 held-out test splits of the MM-BRIGHT benchmark.

\subsection{Evaluation Metrics}
To evaluate the zero-shot retrieval quality, we utilize two standard information retrieval metrics:
\begin{enumerate}
    \item \textbf{Mean Reciprocal Rank (MRR@K):} MRR measures the speed at which the search engine retrieves the first correct (gold) document. It is defined as the average of the reciprocal ranks of the first relevant document across the set of queries $Q$:
    \begin{equation}
    \text{MRR} = \frac{1}{|Q|} \sum_{q \in Q} \frac{1}{\text{rank}_q}
    \end{equation}
    where $\text{rank}_q$ is the position of the first relevant document for query $q$ (assigning a score of $0$ if no relevant document is retrieved in the top $K$ positions).
    
    \item \textbf{Normalized Discounted Cumulative Gain (NDCG@K):} NDCG evaluates the quality of the entire ranked candidate list. It accumulates relevance gains ($\text{rel}_i$) from top to bottom, discounting them logarithmically to model user decay in scanning lower ranks. Discounted Cumulative Gain (DCG) at rank $K$ is defined as:
    \begin{equation}
    \text{DCG}@K = \sum_{i=1}^K \frac{\text{rel}_i}{\log_2(i + 1)}
    \end{equation}
    To make scores comparable across queries, the actual DCG is normalized by the Ideal DCG (IDCG), which represents the perfect rank ordering of all relevant document targets:
    \begin{equation}
    \text{NDCG}@K = \frac{\text{DCG}@K}{\text{IDCG}@K}
    \end{equation}
    Yielding a value in $[0.0, 1.0]$ where $1.0$ indicates perfect ranking.
\end{enumerate}

\subsection{Empirical Evaluation on Held-Out Test Domains}
Following target dataset expansion (scaling to 2,607 targets generated by the \textbf{Teacher Model} \texttt{meta/llama-3.1-70b-instruct} via the NVIDIA NIM API and \texttt{beir} corpora), we evaluated the distilled \textbf{Student Model} (built on the \texttt{Qwen/Qwen2.5-7B-Instruct} backbone with trained Query LUT) across the held-out evaluation splits under the optimal score-fusion regime ($\beta=0.6$) with Pseudo-Relevance Feedback (PRF) enabled ($K=3, \alpha=0.3$).

Table~\ref{tab:new_heldout_results} presents the comparative retrieval metrics (MRR@10 and NDCG@10) across these test domains between the active online \textbf{HIVE Teacher} and our CPU-only \textbf{Distilled Hybrid Student}.

\begin{table}[htbp]
\centering
\scriptsize
\caption{Zero-Shot Cross-Domain Retrieval Metrics (MRR@10 / NDCG@10) on held-out test splits comparing the active online HIVE Teacher (\texttt{meta/llama-3.1-70b-instruct}) against the CPU-only Distilled Hybrid Student (\texttt{Qwen/Qwen2.5-7B-Instruct}, $\beta=0.6$, with PRF).}
\label{tab:new_heldout_results}
\begin{tabular}{lcccc}
\toprule
 & \multicolumn{2}{c}{\textbf{HIVE Teacher (70B/32B)}} & \multicolumn{2}{c}{\textbf{Hybrid Student (7B, Ours)}} \\
\cmidrule(r){2-3} \cmidrule(l){4-5}
\textbf{Held-Out Test Split} & \textbf{MRR@10} & \textbf{NDCG@10} & \textbf{MRR@10} & \textbf{NDCG@10} \\
\midrule
law (30 queries) & 0.5559 & 0.6735 & 0.4000 & 0.4352 \\
gaming (26 queries) & 0.5036 & 0.6106 & \textbf{0.3618} & \textbf{0.4437} \\
askubuntu (35 queries) & 0.3730 & 0.4958 & 0.3027 & 0.3249 \\
salesforce (10 queries) & 0.4225 & 0.5348 & 0.2851 & 0.3352 \\
math (45 queries) & 0.3156 & 0.4254 & 0.1391 & 0.1305 \\
\midrule
\textbf{AVERAGE (OOD Held-Out)} & \textbf{0.4195} & \textbf{0.5337} & \textbf{0.2978} & \textbf{0.3339} \\
\bottomrule
\end{tabular}
\end{table}

\subsection{Teacher-Student Capacity Analysis}
Comparing the active online HIVE Teacher against our distilled Hybrid Student reveals key empirical findings regarding distillation efficiency and domain generalization:
\begin{itemize}
    \item \textbf{Inference-Free Efficiency Retention:} Overall, our CPU-only student pipeline recovers \textbf{71.0\% of the teacher's MRR@10} ($0.2978$ vs $0.4195$) and \textbf{62.6\% of the teacher's NDCG@10} ($0.3339$ vs $0.5337$), while operating in under 5ms on standard CPU hardware without invoking active LLMs or GPUs at query time.
    \item \textbf{High-Efficiency Technical Retention:} On software diagnostic queries (\texttt{askubuntu}), the student retains \textbf{81.2\% of the teacher's MRR@10} ($0.3027$ vs $0.3730$). On legal and gaming splits (\texttt{law} and \texttt{gaming}), the student achieves strong top-10 scores of $0.4352$ and $0.4437$ NDCG@10 respectively.
    \item \textbf{Mathematical Formula Gap:} The largest performance gap remains in \texttt{math}, where the teacher reaches $0.4254$ NDCG@10 while the student reaches $0.1305$. This highlights that complex mathematical derivations require explicit LaTeX formula targets (e.g. ARQMath) during Query LUT training to bridge the formula representation gap.
\end{itemize}

% \section{Research Challenges and Proposed Mitigations}

% \subsection{Encoder Capacity Bottleneck}
% \textbf{Challenge:} Can a visual transformer (ViT) with a vocabulary-projection head match the capacity of a large autoregressive language model? 
% \textbf{Mitigation:} We propose utilizing a deep visual projection model (incorporating late-interaction visual hidden states like ColPali) and optimizing the vocabulary size $|V|$. We also apply high temperatures ($\tau \approx 2.0$) during early training epochs to soften target distributions, allowing the student to learn coarse associations before sharpening lexical weights.

% \subsection{Teacher Generation Cost}
% \textbf{Challenge:} Running HIVE to generate $\hat{q}$ and $R(d, q)$ on massive datasets is financially and computationally expensive.
% \textbf{Mitigation:} We propose an active-learning-based distillation loop. Instead of labeling the entire corpus, we query HIVE only for document pairs that lie close to the decision boundary of an initial baseline V-SPLADE model. This reduces the required LLM API queries by up to 90\% while preserving distillation quality.

\end{document}
